\clearpage\begingroup\NoSourceHead\SourceLeaf{531b}
\renewcommand{\thefootnote}{\arabic{footnote})}\setcounter{footnote}{35}
\ArithmeticHead{74}{L}{BOLZA, GAUSS UND DIE VARIATIONSRECHNUNG.}
\begin{center}\textls[100]{3. Die \textsc{Landsberg}sche Verallgemeinerung des \textsc{Gauss-Bonnet}schen}\\
\textls[100]{Satzes.}\end{center}
\noindent\hspace*{\parindent}18. Der \textsc{Gauss-Bonnet}sche Satz über die Totalkrümmung ist nun selbst
wieder in einem allgemeineren Satz der Variationsrechnung enthalten, wie
\clearpage\SourceLeaf{532}
\ArithmeticHead{75}{R}{DER SATZ ÜBER DIE TOTALKRÜMMUNG.}
\noindent zuerst \textsc{Landsberg}\footnote{G. \textsc{Landsberg}, loc. cit. 29), S. 330—333. Vgl. dazu ferner P. \textsc{Funk} und L. \textsc{Berwald}, Lotos, Prag
67/68 (1920), S. 45.} nachgewiesen hat. Freilich handelt es sich hier \textls[100]{nicht}
um eine Ausdehnung des Satzes auf das \textls[100]{allgemeine} Integral
\[
J=\int_{t_0}^{t_1}F(x,y,x',y')\,dt,
\]
wie bei den früher betrachteten Verallgemeinerungen von Sätzen über geodätische Linien, vielmehr muss die Funktion $F$ sehr starken Einschränkungen
unterworfen werden. Bei der gesuchten Verallgemeinerung entspricht zunächst,
wie bereits oben im Absatz 1, der geodätischen Krümmung die durch (43)
definierte »extremale Krümmung«, dem Längenintegral $s$ das Integral $J$ und
daher dem Differential $ds$ der Integrand $Fdt$, sodass an Stelle der Gleichung
(44) die folgende tritt
\[
K_eFdt=\sqrt{\frac{F_1}{F}}(x'\,dy'-y'\,dx')+\frac{(F_{xy'}-F_{yx'})\,dt}{F_1^{\frac12}F^{\frac12}}.
\]
\par\noindent\hspace*{\parindent}Die Verallgemeinerung des \textsc{Gauss}schen Winkels $\theta$ wird dann eine Funktion
$\Theta(x,y,x',y')$ sein, deren Differential in seinen auf $x'$, $y'$ bezüglichen Gliedern
mit dem ersten Glied des eben hingeschriebenen Ausdrucks übereinstimmt;
d. h. $\Theta$ muss den beiden Bedingungen genügen
\[
\frac{\partial\Theta}{\partial x'}=-y'\sqrt{\frac{F_1}{F}},\quad
\frac{\partial\Theta}{\partial y'}=x'\sqrt{\frac{F_1}{F}}.
\]
Da die Integrabilitätsbedingung erfüllt ist, so ist hierdurch $\Theta$ bis auf eine
additive, willkürlich bleibende Funktion von $x$, $y$ bestimmt. Man erhält jetzt
also als Verallgemeinerung der Gleichung (45)
\begin{equation*}\tag{48}
K_eFdt-d\Theta=S(x,y,x',y')\,dt,
\end{equation*}
wo
\[
S(x,y,x',y')=\frac{F_{xy'}-F_{yx'}}{F_1^{\frac12}F^{\frac12}}-\frac{\partial\Theta}{\partial x}x'-\frac{\partial\Theta}{\partial y}y'.
\]
Hier versagt nun aber der Versuch, die Analogie bei voller Allgemeinheit der
Funktion $F$ weiter zu treiben, da im allgemeinen die Funktion $S(x,y,x',y')$
\textls[100]{nicht} homogen und linear in $x'$, $y'$ ist, wie es bei dem speziellen Problem
der geodätischen Linien der Fall war.
\par\noindent\hfill 10$^{\ast}$\hspace*{5em}\par
\clearpage\SourceLeaf{533}
\ArithmeticHead{76}{L}{BOLZA, GAUSS UND DIE VARIATIONSRECHNUNG.}
\noindent\hspace*{\parindent}In dem speziellen Fall jedoch, wo $S(x,y,x',y')$ homogen und linear in
$x'$, $y'$ ist, gilt in der Tat eine Verallgemeinerung des Satzes (47), wie \textsc{Landsberg} im einzelnen nachweist. Integriert man nämlich die Gleichung (48) in
positivem Sinn um die Begrenzung $\mathfrak C$ eines einfach zusammenhängenden Bereiches $\mathfrak R$ der $x,y$-Ebene, so ergibt sich
\begin{equation*}\tag{49}
\iint\limits_{\mathfrak R}(P_y-Q_x)\,dx\,dy=\int\limits_{\mathfrak C}d\Theta-\int\limits_{\mathfrak C}K_eFdt.
\end{equation*}
\par\noindent\hspace*{\parindent}Die drei hier auftretenden Integrale sind, ebenso wie dies entsprechend
in der Formel (47) der Fall war, absolute Invarianten der Funktion $F$ in
Bezug auf die Gesamtheit aller Punkttransformationen der $x,y$-Ebene und
ihrer Erweiterungen. Ebenso ist das Krümmungsmass $K$ eine absolute Invariante; dagegen ist die Differenz $P_y-Q_x$ keine absolute, sondern nur eine
Invariante vom Index\footnote{Wegen der Terminologie vergleiche A. L. \textsc{Underhill}, loc. cit. 30), S. 320.} 1. Wenn daher eine zweite Invariante $\mathfrak I$ von $F$
vom Index 1 existiert, so lässt sich eine vollständige Analogie mit dem Fall
der geodätischen Linien herstellen, indem man den Quotienten
\begin{equation*}\tag{50}
K=\frac{P_y-Q_x}{\mathfrak I},
\end{equation*}
der nunmehr ebenfalls eine absolute Invariante ist, als Verallgemeinerung des
Krümmungsmasses und gleichzeitig das Produkt $\mathfrak I\,dx\,dy$ als Verallgemeinerung
des Flächenelementes $d\sigma$ auffasst. Die Gleichung (49) erscheint dann in der
Tat als eine naturgemässe Verallgemeinerung des \textsc{Gauss-Bonnet}schen Satzes (47).
\par\begin{center}\rule{.18\linewidth}{.3pt}\end{center}
\clearpage\endgroup
